% =============================================================================
%  Capítulo 2 — Uso Básico do Hanaro (Parte II, Passo a Passo)
%  Capítulo único da parte, só em seções: os oito passos, as colunas do
%  histórico de execuções, as telas do Modo TV e o encerramento da sessão.
%
%  Cada \passo é uma linha de duas colunas: a tela à esquerda e o texto à
%  direita, e conta como uma seção do capítulo (2.1, 2.2...). O argumento
%  opcional lista as posições dos marcadores sobre a tela, como x/y em fração
%  da imagem, com origem no canto SUPERIOR esquerdo. A numeração segue a
%  ordem da lista e casa com a lista `elementos'.
%
%      \passo[0.335/0.525, 0.665/0.649]{01-login}{Título}{Texto ...}
% =============================================================================
\chapter{Uso Básico do Hanaro}\label{cap:passos}

% ---------- acesso e atualização da base ----------
\passo[0.335/0.525, 0.665/0.649, 0.335/0.772]{01-login}{Entrar no Hanaro}{%
    Abra o endereço do Hanaro no navegador. A tela de acesso é a primeira a aparecer.
    \begin{elementos}
        \item \campo{Usuário}: o nome de usuário cadastrado pelo administrador. Não é
              o e-mail.
        \item \campo{Senha}: a senha do usuário. O ícone de olho, no fim do campo,
              mostra o que foi digitado.
        \item \botao{Entrar}: valida as credenciais e abre o Dashboard.
    \end{elementos}
    Credenciais inválidas mantêm o usuário nesta tela.
}

\passo[0.115/0.181, 0.845/0.146, 0.31/0.521, 0.69/0.669]{02-enviar-relatorio}{Enviar o relatório do GERP}{%
    A base de scrap é atualizada pelo relatório \emph{Other Account Transaction
    Text}. Exporte-o do GERP antes deste passo.
    \begin{elementos}
        \item \caminho{Execuções}, no menu lateral: abre a tela
              \textbf{Execuções e atualização}.
        \item \botao{Enviar relatório}: abre a janela \textbf{Atualizar base por
              relatório}.
        \item \textbf{Área do arquivo}: um clique abre a escolha do arquivo exportado
              do GERP. O nome do arquivo escolhido aparece no lugar de
              \emph{Selecionar relatório GERP}.
        \item \botao{Processar relatório}: envia o arquivo. O Hanaro registra uma
              nova execução, conferida no passo seguinte.
    \end{elementos}
    \botao{Cancelar} fecha a janela sem enviar.
}

\begin{remark}
O arquivo precisa ser o \emph{Other Account Transaction Text}, em texto separado por
tabulação (TSV), com codificação UTF-8 ou CP1252, como indica a própria janela.
Planilha salva em outro formato não serve.
\end{remark}

\passo[0.148/0.378, 0.98/0.378, 0.871/0.242]{03-execucoes}{Conferir a execução}{%
    A execução nova entra no topo do \textbf{Histórico de rotinas}.
    \begin{elementos}
        \item \textbf{Linha da execução}: o identificador (\texttt{EXE-...}), a
              origem e o gatilho. Um clique na linha abre os \textbf{Detalhes da
              execução}, com o arquivo de origem e as contagens.
        \item \textbf{Status}: \emph{Na fila} enquanto a execução aguarda o
              processamento e \emph{Concluído} ao final. Ao lado, o
              \textbf{Snapshot} passa de \emph{Não publicado} a \emph{Publicado}.
        \item \botao{Filtrar}: restringe o histórico.
    \end{elementos}
    As colunas estão descritas na \autoref{tab:historico}.
}

\begin{remark}[info]
O Dashboard passa a refletir o arquivo novo quando o snapshot da execução fica
\emph{Publicado}. Até lá, ele mostra a base da execução anterior.
\end{remark}

\section{Colunas do histórico}

% Tabela curta: vai inteira, em `table', para não partir entre páginas.
\begin{table}[H]
\caption{Colunas do histórico de rotinas.}\label{tab:historico}
\guiazebra
\small
\noindent\begin{tabularx}{\textwidth}{@{}L{4.4cm}X@{}}
    \toprule
    \guiacabtabela
    \thead{Coluna} & \thead{Conteúdo} \\
    \midrule
    \textbf{Execution ID / Processo} & Identificador da execução e nome do relatório
                                       processado \\
    \textbf{Origem}                  & Sistema de onde veio o arquivo: GERP \\
    \textbf{Gatilho}                 & Como a execução começou. \emph{Manual}: pelo
                                       envio do relatório no passo~2 \\
    \textbf{Início}, \textbf{Fim} e \textbf{Duração}
                                     & Data e hora do processamento e o tempo que ele
                                       levou \\
    \textbf{Recebidos}               & Registros lidos no arquivo \\
    \textbf{Válidos}                 & Registros aceitos na base de scrap \\
    \textbf{Rejeitados}              & Registros recusados na validação \\
    \textbf{Snapshot}                & \emph{Publicado} quando a base processada já
                                       alimenta os indicadores; \emph{Não publicado}
                                       antes disso \\
    \textbf{Status}                  & \emph{Na fila} enquanto aguarda;
                                       \emph{Concluído} ao final \\
    \bottomrule
\end{tabularx}
\end{table}

\begin{remark}
Com \textbf{Rejeitados} acima de zero, parte do arquivo ficou fora da base. Abra os
detalhes da execução e acione o suporte em \caminho{Ajuda e suporte}, no pé do menu
lateral, com o identificador da execução.
\end{remark}

% ---------- indicadores e ocorrências ----------
\passo[0.30/0.114, 0.36/0.214, 0.66/0.114, 0.45/0.322, 0.924/0.503, 0.27/0.725]{04-dashboard}{Ler o Dashboard}{%
    O Dashboard é a tela que abre depois do acesso. No alto, à direita,
    \emph{Atualizado às} indica a hora da última atualização.
    \begin{elementos}
        \item \textbf{Impacto absoluto}: chave que alterna os indicadores entre o IF
              Cost em dólar e o Scrap Rate, em percentual.
        \item \textbf{Indicadores}: IF Cost acumulado, target, atingimento do target,
              diferença para a meta e variação sobre o mesmo acumulado do ano
              anterior.
        \item \textbf{Controles}: o botão com o ícone de tela abre o Modo TV
              (passo~8); \botao{YTD anterior} escolhe a base de comparação;
              \botao{Filtros} recorta os dados; \textbf{Valor monetário} alterna
              entre dólar e quantidade, em unidades.
        \item \textbf{Target × realizado mensal}: realizado do ano, referência do ano
              anterior e target, mês a mês.
        \item \textbf{Top 10 produtos / áreas}: o IF Cost por produto ou área.
        \item \textbf{Top 10 linhas} e \textbf{Top 10 modelos}: os rankings de
              ofensores. Mais abaixo, componentes e ofensores.
    \end{elementos}
}

\begin{remark}[dica]
Cada quadro tem o seu botão \botao{Filtros} e um botão de tabela, que abre os mesmos
números em linhas e colunas. Passar o cursor sobre um ponto do gráfico mostra os
valores do mês.
\end{remark}

\passo[0.12/0.178, 0.555/0.194, 0.135/0.337, 0.70/0.1]{05-base-scrap}{Revisar a Base de Scrap}{%
    A lista traz as ocorrências consolidadas pelas execuções. Ao lado do título, o
    total de ocorrências, o valor em dólar e a quantidade.
    \begin{elementos}
        \item \caminho{Base de Scrap}, no menu lateral: abre a tela.
        \item \textbf{Ordem e filtros}: \botao{Data da transação} escolhe o campo da
              ordem e \botao{Decrescente} a inverte; a lista \botao{Todos} filtra
              pelo status da análise; \botao{Filtrar} abre os demais filtros.
        \item \textbf{Ocorrência}: organização e produto, data, item, descrição,
              ordem, quantidade, valor, tipo de scrap, status da análise
              (\emph{Não analisado} ou \emph{Revisado}), responsável e relatório. Um
              clique abre a análise (passo~6).
        \item \botao{Aplicar Modelo} e \botao{Selecionar itens}: levam uma análise já
              feita a várias ocorrências (passo~7).
    \end{elementos}
}

\passo[0.39/0.138, 0.195/0.349, 0.195/0.507, 0.195/0.652, 0.577/0.87]{06-ocorrencia}{Consultar a análise de uma ocorrência}{%
    A janela \textbf{Análise de scrap} reúne o que se sabe da ocorrência. Role a
    janela para ver todas as seções.
    \begin{elementos}
        \item \textbf{Status da análise}: \emph{Revisado}, como nesta ocorrência, ou
              \emph{Não analisado}.
        \item \textbf{Custo}: tipo de custo, modelo e quantidade. Acima, \textbf{Onde
              e quando} traz a data e a hora.
        \item \textbf{Contenção imediata}: a ação tomada na linha para conter o
              problema.
        \item \textbf{Evidências e fotos}: as imagens anexadas à análise.
        \item \botao{Criar relatório desta análise}: abre um relatório de análise a
              partir desta ocorrência.
    \end{elementos}
    O coração, no alto, guarda a análise como modelo para o passo seguinte.
}

\passo[0.565/0.09, 0.145/0.494, 0.313/0.632, 0.685/0.768]{07-modelo-lote}{Aplicar um modelo em lote}{%
    Quando o mesmo problema se repete em várias ocorrências, a análise de uma serve
    de modelo para as outras.
    \begin{elementos}
        \item \textbf{Modelo em uso}: em \botao{Aplicar Modelo}, escolha o modelo em
              \textbf{Modelos salvos} e clique em \botao{Usar este modelo}. A faixa
              no alto confirma a escolha.
        \item \textbf{Ocorrências marcadas}: com \botao{Selecionar itens}, marque as
              ocorrências e clique em \botao{Aplicar modelo aos itens selecionados},
              acima da lista. Abre esta janela.
        \item \textbf{Reaproveitar imagens do relatório de referência}: copia as
              fotos do modelo para todas as ocorrências marcadas.
        \item \botao{Aplicar modelo aos itens selecionados}: grava a análise do
              modelo nas ocorrências marcadas.
    \end{elementos}
}

\begin{remark}
Ocorrências inativas ou que já têm relatório são ignoradas na aplicação, como avisa a
própria janela. Copiar as fotos replica as mesmas evidências em todas as ocorrências
marcadas: use só quando as fotos valem para todas.
\end{remark}

% ---------- Modo TV ----------
\passo[0.395/0.069, 0.712/0.083, 0.2/0.865, 0.5/0.865, 0.785/0.865]{08-modo-tv}{Exibir o Modo TV}{%
    O Modo TV mostra os indicadores em tela cheia, para as TVs da fábrica. Ele abre
    pelo botão com o ícone de tela, nos controles do Dashboard (passo~4).
    \begin{elementos}
        \item \textbf{Tela em exibição}: o título da tela. Na tela Fábrica,
              \botao{Mês} e \botao{Ano} escolhem o período.
        \item \textbf{Dados até}: a data do último dado da base e, ao lado, a data e
              a hora atuais.
        \item \textbf{Telas}: \botao{Executivo}, \botao{Fábrica}, \botao{Ofensores} e
              \botao{Prioridades}, descritas na \autoref{tab:telastv}. Um clique leva
              direto à tela.
        \item \textbf{Rotação}: as setas passam de tela, \botao{Pausar} para a
              rotação automática e \textbf{Intervalo} define o tempo de cada tela.
        \item \botao{Tema Claro}, \botao{Tela cheia} e \botao{Sair}: trocam o fundo
              escuro pelo claro, ocupam a tela inteira e voltam ao Dashboard.
    \end{elementos}
}

\begin{remark}[dica]
Na TV da linha, deixe a \textbf{Rotação automática} ligada e use \botao{Tela cheia}:
as quatro telas se alternam sem intervenção. O contador no fim da barra, como
\emph{2/4}, indica a tela atual.
\end{remark}

\section{Telas do Modo TV}

% Tabela curta: vai inteira, em `table', para não partir entre páginas.
\begin{table}[H]
\caption{As quatro telas do Modo TV.}\label{tab:telastv}
\guiazebra
\small
\noindent\begin{tabularx}{\textwidth}{@{}L{2.4cm}L{3.9cm}X@{}}
    \toprule
    \guiacabtabela
    \thead{Botão} & \thead{Título da tela} & \thead{O que mostra} \\
    \midrule
    \textbf{Executivo}   & Visão Executiva
                         & IF Cost acumulado, mesmo período do ano anterior, variação e
                           target acumulado; evolução do IF Cost no ano, no ano anterior
                           e no target; os três componentes mais afetados \\
    \textbf{Fábrica}     & Ranking de ocorrências por linha
                         & Ocorrências do período, linha com mais ocorrências, linhas
                           monitoradas e período anterior; ranking das linhas de
                           montagem, com a variação; setores com mais ocorrências \\
    \textbf{Ofensores}   & Principais ofensores do período
                         & Os cinco defeitos e as oito linhas de maior IF Cost; o part
                           number e o defeito mais críticos \\
    \textbf{Prioridades} & Ocorrências prioritárias
                         & As ocorrências críticas de maior custo, com a descrição;
                           situação das revisões (pendentes, em revisão e
                           justificadas) e categorias predominantes \\
    \bottomrule
\end{tabularx}
\end{table}

% --- Passos complementares adicionados a partir das capturas reais ---
\passo[0.065/0.415, 0.280/0.205, 0.930/0.175]{26-relatorios-lista}{Consultar e criar Relatórios}{%
  A tela \textbf{Relatórios} reúne os planos de ação abertos para tratar as ocorrências analisadas na Base de Scrap.
  \begin{elementos}
    \item \caminho{Relatórios}, no menu lateral: abre a central de relatórios e planos de ação.
    \item \textbf{Abas de status}: filtram rapidamente entre \botao{Todos}, \botao{Abertos}, \botao{Em andamento} e \botao{Fechados}.
    \item \botao{+ Novo relatório}: inicia o assistente de 4 etapas (vinculação de análises revisadas, dados do relatório, 5 porquês e ações iniciais).
  \end{elementos}
}

\passo[0.320/0.195, 0.610/0.465, 0.935/0.145]{34-relatorio-tarefas-quadro-validar}{Acompanhar e validar tarefas no Quadro}{%
  Dentro de um relatório, a aba \textbf{Quadro} organiza as tarefas do plano de ação em quatro colunas: \emph{Planejada}, \emph{Em andamento}, \emph{Em validação} e \emph{Concluída}.
  \begin{elementos}
    \item \botao{Quadro}: alterna a visualização do plano de ação para o formato de colunas.
    \item \botao{Validar conclusão}: disponível na coluna \emph{Em validação} para que o gestor aprove a entrega da tarefa e a mova para \emph{Concluída}.
    \item \botao{+ Nova tarefa}: abre o formulário completo para cadastrar ação, responsável, área, prazo e economia prevista.
  \end{elementos}
}

\passo[0.320/0.468, 0.560/0.228, 0.718/0.838]{40-relatorio-exportar-pdf}{Exportar documento em PDF ou Apresentação}{%
  O botão \botao{Baixar documento}, no topo do relatório, monta o arquivo executivo para leitura ou reunião gerencial.
  \begin{elementos}
    \item \textbf{Formato e idioma}: escolha entre \botao{PDF A4} e \botao{Apresentação 16:9} (PowerPoint) e selecione o idioma (\emph{Português}, \emph{English} ou \emph{Coreano}).
    \item \campo{Traduzir o texto para inglês (IA)}: traduz automaticamente os textos abertos antes do download.
    \item \botao{Baixar PDF} (ou \botao{Baixar apresentação}): gera o documento com as seções marcadas.
  \end{elementos}
}

\passo[0.065/0.440, 0.930/0.175, 0.930/0.355]{42-alertas-historico}{Acompanhar e configurar Alertas}{%
  A tela \textbf{Alertas} exibe as notificações disparadas por tarefas atribuídas, prazos vencidos ou regras de limite de custo.
  \begin{elementos}
    \item \caminho{Alertas}, no menu lateral: abre o \textbf{Histórico de notificações}.
    \item \botao{Configurar alertas}: abre a janela de regras de gatilho por valor, meta, lembrete ou evento.
    \item \textbf{Ações da linha}: \botao{Ver tarefa} abre o item correspondente e o ícone de confirmação marca o aviso como lido.
  \end{elementos}
}

\section{Encerramento da sessão}

O nome do usuário fica no pé do menu lateral. O botão \botao{\faEllipsisV}, ao lado
dele, leva ao \textbf{Perfil}. Ali, a seção \textbf{Sessões ativas} lista os
dispositivos conectados à conta: \botao{Encerrar sessão} desconecta um deles e
\botao{Encerrar outras sessões} desconecta todos, menos o atual.

Em computador compartilhado, encerre a sessão ao terminar e feche o navegador.

\fimdecapitulo
